\documentclass[11pt]{article}
\newcommand{\ListaTitulo}{Gabarito 07: Teorema de Bayes}
% Preâmbulo compartilhado: MATF14, Listas de Exercícios 2026
% Usado por ListaNN.tex e GabaritoNN.tex via % Preâmbulo compartilhado: MATF14, Listas de Exercícios 2026
% Usado por ListaNN.tex e GabaritoNN.tex via % Preâmbulo compartilhado: MATF14, Listas de Exercícios 2026
% Usado por ListaNN.tex e GabaritoNN.tex via \input{preamble}

\usepackage[utf8]{inputenc}
\usepackage[T1]{fontenc}
\usepackage[brazil]{babel}
\usepackage{fullpage}
\usepackage{amsmath,amssymb}
\usepackage{enumitem}
\usepackage{xcolor}
\usepackage{fancyhdr}
\usepackage{titlesec}
\usepackage{listings}
\usepackage{hyperref}
\usepackage{booktabs}
\usepackage{graphicx}
\usepackage{tikz}

\definecolor{matf14azul}{HTML}{0D3B54}
\definecolor{matf14dourado}{HTML}{C9971E}

\titleformat{\section}{\normalfont\Large\bfseries\color{matf14azul}}{\thesection}{1em}{}
\titleformat{\subsection}{\normalfont\large\bfseries\color{matf14azul}}{\thesubsection}{1em}{}

\providecommand{\ListaTitulo}{Lista de Exercícios}

\setlength{\headheight}{14pt}
\pagestyle{fancy}
\fancyhf{}
\lhead{\small MATF14: Estatística Econômica I}
\rhead{\small \ListaTitulo}
\cfoot{\thepage}
\renewcommand{\headrulewidth}{0.4pt}

\lstset{
  basicstyle=\ttfamily\small,
  breaklines=true,
  frame=single,
  columns=fullflexible,
  backgroundcolor=\color{gray!8},
  commentstyle=\color{matf14dourado},
  keywordstyle=\color{matf14azul}\bfseries,
  language=R
}

\newcounter{questaocounter}
\newenvironment{questao}[1]{%
  \stepcounter{questaocounter}%
  \bigskip\noindent\textbf{Questão #1.}\ %
}{\par}

\newcommand{\cabecalho}[2]{%
  \begin{center}
    {\Large\bfseries\color{matf14azul} #1}\\[0.3em]
    {\large #2}
  \end{center}
  \vspace{0.5em}
  \noindent\rule{\linewidth}{0.6pt}
  \vspace{0.5em}
}

\newenvironment{solucao}{%
  \par\smallskip\noindent\textit{\color{matf14dourado!80!black}Solução.}\ %
}{\hfill$\blacksquare$\par}


\usepackage[utf8]{inputenc}
\usepackage[T1]{fontenc}
\usepackage[brazil]{babel}
\usepackage{fullpage}
\usepackage{amsmath,amssymb}
\usepackage{enumitem}
\usepackage{xcolor}
\usepackage{fancyhdr}
\usepackage{titlesec}
\usepackage{listings}
\usepackage{hyperref}
\usepackage{booktabs}
\usepackage{graphicx}
\usepackage{tikz}

\definecolor{matf14azul}{HTML}{0D3B54}
\definecolor{matf14dourado}{HTML}{C9971E}

\titleformat{\section}{\normalfont\Large\bfseries\color{matf14azul}}{\thesection}{1em}{}
\titleformat{\subsection}{\normalfont\large\bfseries\color{matf14azul}}{\thesubsection}{1em}{}

\providecommand{\ListaTitulo}{Lista de Exercícios}

\setlength{\headheight}{14pt}
\pagestyle{fancy}
\fancyhf{}
\lhead{\small MATF14: Estatística Econômica I}
\rhead{\small \ListaTitulo}
\cfoot{\thepage}
\renewcommand{\headrulewidth}{0.4pt}

\lstset{
  basicstyle=\ttfamily\small,
  breaklines=true,
  frame=single,
  columns=fullflexible,
  backgroundcolor=\color{gray!8},
  commentstyle=\color{matf14dourado},
  keywordstyle=\color{matf14azul}\bfseries,
  language=R
}

\newcounter{questaocounter}
\newenvironment{questao}[1]{%
  \stepcounter{questaocounter}%
  \bigskip\noindent\textbf{Questão #1.}\ %
}{\par}

\newcommand{\cabecalho}[2]{%
  \begin{center}
    {\Large\bfseries\color{matf14azul} #1}\\[0.3em]
    {\large #2}
  \end{center}
  \vspace{0.5em}
  \noindent\rule{\linewidth}{0.6pt}
  \vspace{0.5em}
}

\newenvironment{solucao}{%
  \par\smallskip\noindent\textit{\color{matf14dourado!80!black}Solução.}\ %
}{\hfill$\blacksquare$\par}


\usepackage[utf8]{inputenc}
\usepackage[T1]{fontenc}
\usepackage[brazil]{babel}
\usepackage{fullpage}
\usepackage{amsmath,amssymb}
\usepackage{enumitem}
\usepackage{xcolor}
\usepackage{fancyhdr}
\usepackage{titlesec}
\usepackage{listings}
\usepackage{hyperref}
\usepackage{booktabs}
\usepackage{graphicx}
\usepackage{tikz}

\definecolor{matf14azul}{HTML}{0D3B54}
\definecolor{matf14dourado}{HTML}{C9971E}

\titleformat{\section}{\normalfont\Large\bfseries\color{matf14azul}}{\thesection}{1em}{}
\titleformat{\subsection}{\normalfont\large\bfseries\color{matf14azul}}{\thesubsection}{1em}{}

\providecommand{\ListaTitulo}{Lista de Exercícios}

\setlength{\headheight}{14pt}
\pagestyle{fancy}
\fancyhf{}
\lhead{\small MATF14: Estatística Econômica I}
\rhead{\small \ListaTitulo}
\cfoot{\thepage}
\renewcommand{\headrulewidth}{0.4pt}

\lstset{
  basicstyle=\ttfamily\small,
  breaklines=true,
  frame=single,
  columns=fullflexible,
  backgroundcolor=\color{gray!8},
  commentstyle=\color{matf14dourado},
  keywordstyle=\color{matf14azul}\bfseries,
  language=R
}

\newcounter{questaocounter}
\newenvironment{questao}[1]{%
  \stepcounter{questaocounter}%
  \bigskip\noindent\textbf{Questão #1.}\ %
}{\par}

\newcommand{\cabecalho}[2]{%
  \begin{center}
    {\Large\bfseries\color{matf14azul} #1}\\[0.3em]
    {\large #2}
  \end{center}
  \vspace{0.5em}
  \noindent\rule{\linewidth}{0.6pt}
  \vspace{0.5em}
}

\newenvironment{solucao}{%
  \par\smallskip\noindent\textit{\color{matf14dourado!80!black}Solução.}\ %
}{\hfill$\blacksquare$\par}


\begin{document}

\cabecalho{Gabarito 07}{Teorema da probabilidade total e Teorema de Bayes (Cap. 2, Aula 14)}

\begin{questao}{1} \textbf{(Probabilidade total: inadimplência em carteira segmentada)}
\begin{solucao}
\begin{enumerate}[label=(\alph*)]
  \item $P(\text{inad}) = 0{,}6\times 0{,}04 + 0{,}3\times 0{,}07 + 0{,}1\times 0{,}01
    = 0{,}024+0{,}021+0{,}001 = 0{,}046$ (4,6\%).
  \item Não necessariamente. A inadimplência geral observada (6\%) pode subir mesmo com taxas
    por segmento inalteradas, se a \textbf{composição} da carteira mudou no trimestre, por
    exemplo, se a participação de Pequena Empresa (segmento de maior taxa, 7\%) cresceu às custas
    de Grande Empresa (menor taxa, 1\%), a média ponderada sobe mesmo sem nenhuma taxa individual
    ter piorado.
\end{enumerate}
\end{solucao}
\end{questao}

\begin{questao}{2} \textbf{(Bayes: origem de um sinistro de seguro)}
\begin{solucao}
\begin{enumerate}[label=(\alph*)]
  \item $P(S) = 0{,}75\times 0{,}05 + 0{,}25\times 0{,}20 = 0{,}0375+0{,}05 = 0{,}0875$ (8,75\%).
  \item $P(B\mid S) = \dfrac{P(S\mid B)P(B)}{P(S)} = \dfrac{0{,}20\times 0{,}25}{0{,}0875} =
    \dfrac{0{,}05}{0{,}0875} \approx 0{,}571$ (57,1\%).
  \item A proporção original de Perfil B na carteira é 25\%; depois de observar que houve
    sinistro, a probabilidade de ser Perfil B sobe para 57,1\%, mais que o dobro. Observar
    "houve sinistro" é uma evidência forte a favor do Perfil B (motoristas novatos), porque a taxa
    de sinistro desse perfil (20\%) é bem mais alta que a do Perfil A (5\%): a atualização
    bayesiana desloca a crença fortemente na direção do grupo com maior probabilidade de gerar a
    evidência observada.
\end{enumerate}
\end{solucao}
\end{questao}

\begin{questao}{3} \textbf{(Bayes com três categorias: fraude em declarações de imposto de renda)}
\begin{solucao}
\begin{enumerate}[label=(\alph*)]
  \item $P(\text{sinal}) = 0{,}85\times 0{,}02 + 0{,}12\times 0{,}30 + 0{,}03\times 0{,}90
    = 0{,}017 + 0{,}036 + 0{,}027 = 0{,}080$ (8,0\%).
  \item $P(\text{Normal}\mid\text{sinal}) = 0{,}017/0{,}080 = 0{,}2125$ (21,25\%).
    $P(\text{Atípica}\mid\text{sinal}) = 0{,}036/0{,}080 = 0{,}45$ (45\%).
    $P(\text{Fraude}\mid\text{sinal}) = 0{,}027/0{,}080 = 0{,}3375$ (33,75\%).
  \item Sim, a soma deve dar exatamente 1: $0{,}2125+0{,}45+0{,}3375=1{,}00$. As três categorias
    (Normal, Atípica-legítima, Fraudulenta) formam uma \textbf{partição} do espaço amostral, toda
    declaração sinalizada pertence a exatamente uma delas, e as probabilidades a posteriori de
    uma partição, condicionadas ao mesmo evento, sempre somam 1 (é a mesma lógica do teorema da
    probabilidade total, aplicada "de trás para frente").
\end{enumerate}
\end{solucao}
\end{questao}

\begin{questao}{4} \textbf{(Discussão: base rate e política pública)}
\begin{solucao}
\begin{enumerate}[label=(\alph*)]
  \item $P(\text{doente})=0{,}005$, $P(+\mid\text{doente})=0{,}98$,
    $P(+\mid\text{sadio})=1-0{,}90=0{,}10$.
    $$
    P(+) = 0{,}98\times 0{,}005 + 0{,}10\times 0{,}995 = 0{,}0049+0{,}0995 = 0{,}1044.
    $$
    $$
    P(\text{doente}\mid +) = \frac{0{,}0049}{0{,}1044} \approx 0{,}047 \ (4{,}7\%).
    $$
  \item Mesmo com sensibilidade alta (98\%), apenas cerca de 4,7\% das pessoas com teste positivo
    realmente têm a doença, porque a doença é rara (0,5\%) e a especificidade (90\%) ainda deixa
    passar 10\% de falsos positivos entre a enorme maioria saudável da população. Tratar
    imediatamente todo positivo em uma triagem de massa significaria submeter mais de 95\% dos
    "positivos" (pessoas saudáveis) a um tratamento desnecessário, a proposta ignora a taxa base
    (\emph{base rate}) da doença; a decisão correta seria usar o teste positivo como gatilho para
    um segundo exame confirmatório mais específico, não para tratamento direto.
\end{enumerate}
\end{solucao}
\end{questao}

\end{document}
